\subsection{The completion of a locally convex space}

THEOREM 1 (Grothendieck). --- Let E be a locally convex space, and let S be an adapted and covering bornology on E. Let \(F \subset E^{*}\) be the space of those linear forms on E whose restriction to each set belonging to S is continuous. If F is assigned the S-topology, then the canonical injection from \(E'_{S}\) into F extends to an isomorphism from the completion \(\hat{E}'_{S}\) of \(E'_{S}\) onto F.

Since every simple limit of linear forms on E is a linear form (III, p.~16, prop. 4) and since the bornology \(\mathfrak{S}\) on E is covering, it follows from GT, X, § 1, No.~6, cor. 2 that the space F with the \(\mathfrak{S}\) -topology is Hausdorff and complete. It is clear that \(E_{\mathfrak{S}}'\) is a topological vector subspace of F; hence it is enough to prove that \(E_{\mathfrak{S}}'\) is everywhere dense in F. This follows from the following lemma:

Lemma 1. --- Let A be a closed convex balanced subset of E and let u be a linear form on E whose restriction to A is continuous. Then for every \(\varepsilon > 0\) , there exists an \(x' \in E'\) such that

\[
\left| u (x) - \langle x, x ^ {\prime} \rangle \right| \leqslant \varepsilon \quad \text { for   every } \quad x \in A.
\]

Let \(\varepsilon > 0\) . There exists a closed convex balanced neighbourhood U of 0 in E such that \(|u(x)| \leqslant \varepsilon\) for all \(x \in \mathrm{U} \cap \mathrm{A}\) . We know that the polar \(\mathrm{U}^{\circ}\) of U in \(\mathrm{E}^{*}\) is contained in \(\mathrm{E}'\) and is compact for the topology \(\sigma(\mathrm{E}^{*}, \mathrm{E})\) (III, p.~17, cor. 2). Since the polar \(\mathrm{A}^{\circ}\) of A in \(\mathrm{E}^{*}\) is closed for \(\sigma(\mathrm{E}^{*}, \mathrm{E})\) , it follows that \(\mathrm{A}^{\circ} + \mathrm{U}^{\circ}\) is a closed convex subset of \(\mathrm{E}^{*}\) (GT, III, § 4, No.~1, cor. 1).

Let C be a closed convex balanced subset of E. Then C is closed for \(\sigma(\mathrm{E},\mathrm{E}^{\prime})\) (II, p.~45, cor. 3), hence also for \(\sigma(\mathrm{E},\mathrm{E}^{*})\) , and as a consequence, we have \(C = C^{\circ\circ}\) (for the duality between E and \(E^{*}\) ). As a result, we have

\[
\mathrm{A} \cap \mathrm{U} = \mathrm{A} ^ {\circ \circ} \cap \mathrm{U} ^ {\circ \circ} = (\mathrm{A} ^ {\circ} \cup \mathrm{U} ^ {\circ}) ^ {\circ} \supset (\mathrm{A} ^ {\circ} + \mathrm{U} ^ {\circ}) ^ {\circ}
\]

from which, we get

\[
(\mathrm{A} \cap \mathrm{U}) ^ {\circ} \subset (\mathrm{A} ^ {\circ} + \mathrm{U} ^ {\circ}) ^ {\circ \circ} = \mathrm{A} ^ {\circ} + \mathrm{U} ^ {\circ}.
\]

Since the linear form \(\varepsilon^{-1}u\) belongs to \((\mathrm{A}\cap\mathrm{U})^{\circ}\) , there exist \(v\in A^{\circ}\) and \(w\in U^{\circ}\) such that \(u=\varepsilon(v+w)\) . Hence \(x'=\varepsilon w\) belongs to \(E'\) and \(u-x'=\varepsilon v\) is bounded above in absolute value by \(\varepsilon\) on A; hence the lemma.

Now let E be a locally convex Hausdorff space and \(\hat{E}\) its completion. Every continuous linear form f on E extends to \(\hat{E}\) by continuity: hence we have \((\hat{\mathrm{E}})' = \mathrm{E}'\) (III, p.~16) and every element of \(\hat{\mathbf{E}}\) defines a linear form on \(\mathbf{E}'\) ; that is, an element of the algebraic dual \(\mathbf{E}^{\prime *}\) of \(\mathbf{E}'\) . In addition, the duality between \(\mathbf{E}\) (resp. \(\hat{\mathbf{E}}\) ) and \(\mathbf{E}'\) is separating (II, p.~24, cor. 1). Consequently \(\mathbf{E}\) and \(\hat{\mathbf{E}}\) can be identified with vector subspaces of \(\mathbf{E}^{\prime *}\) .

THEOREM 2. --- Let E be a locally convex Hausdorff space and \(\hat{\mathbf{E}}\) its completion; we identify E and \(\hat{\mathbf{E}}\) with vector subspaces of \(\mathbf{E}^{\prime *}\) . Then for an element \(f \in \mathbf{E}^{\prime *}\) to belong to \(\hat{\mathbf{E}}\) , it is necessary and sufficient that the restriction of \(f\) to every equicontinuous subset of \(\mathbf{E}'\) is continuous for the topology \(\sigma(\mathbf{E}', \mathbf{E})\) .

The space E can be identified with the topological dual of \(E'\) when \(E'\) is assigned the topology \(\sigma(E', E)\) (II, p.~43, prop. 3); on the other hand, if S is the set of equicontinuous subsets of \(E'\) , the given topology on E is the S-topology (III, p.~19, cor. 1). Then it follows from III, p.~13, prop. 1, that the sets of S are bounded for \(\sigma(E', E)\) (cf.~later on, III, p.~22, prop. 9); in other words, S is an adapted and covering bornology for the topology \(\sigma(E', E)\) . Theorem 2 is then a consequence of th. 1 if we replace E by \(E'\) and \(E'_{S}\) by E.

COROLLARY 1 (Banach). --- Let E be a Hausdorff and complete locally convex space. In order that a linear form on \(E'\) be continuous for the weak topology \(\sigma(E', E)\) (i.e.~arises from an element of E) it is sufficient that its restriction to every equicontinuous subset of \(E'\) is continuous for \(\sigma(E', E)\) .

Remark. --- Suppose in addition, that there exists a countable total set in E; then every equicontinuous subset of \(E'\) is metrizable for the topology \(\sigma(E', E)\) (III, p.~18, prop. 6); therefore to verify that a linear form u on \(E'\) is weakly continuous, it is enough to verify that for every equicontinuous sequence \((x_{n}')\) in \(E'\) which converges to 0 for \(\sigma(E', E)\) , we have \(\lim_{n \to \infty} u(x_{n}') = 0\) .

COROLLARY 2. --- Let \((\mathrm{E}_{i})_{i\in\mathbb{I}}\) be a family of Hausdorff locally convex spaces and let E be their topological direct sum. Then the canonical mapping from the direct sum of the \(\hat{E}_{i}\) into \(\hat{E}\) is an isomorphism. In particular, E is complete if and only if all the \(E_{i}\) are complete.

We know that the dual of E can be identified with the product of the duals of the \(E_{i}\) (II, p.~30, formula (1)). Let \(u \in \hat{E}\) , and let \(u_{i} \in E_{i}^{\prime *}\) be the restriction of u (considered as an element of \(E^{\prime *}\) ) to \(E_{i}^{\prime} \subset E^{\prime}\) . It is immediate that it is enough to prove that \(u_{i} = 0\) except for a finite number of indices \(i \in I\) . Suppose on the contrary that there exists a sequence \((i_{n})_{n \in \mathbf{N}}\) of distinct indices such that \(u_{i_{n}} \neq 0\) . Then there exists \(x_{i_{n}} \in E_{i_{n}}^{\prime}\) such that \(u_{i_{n}}(x_{i_{n}}) = n\) . The set H of all \(x_{i_{n}}\) is equicontinuous in \(E^{\prime}\) and the restriction of u to H is not bounded, which is impossible.

